% Imporditud: tools/import_eksperimendid.py
% Allikas: Eesti füüsikaolümpiaadi eksperimendi ülesannete
%   kogu (Rael Kalda, 2023), ülesanne Ü110, lahendus L110.
\setAuthor{Erkki Tempel}
\setRound{piirkonnavoor}
\setYear{2015}
\setNumber{G E1}
\setDifficulty{3}
\setTopic{varia}
\setVahendid{Kaks tikku}
\setPunktid{12}
\setAllikas{Eksperimendi kogu Ü110, lahendus lk 85}
\setLahendusseis{toores_ilma_skeemita}

\prob{Tikud}
Mitu protsenti moodustab tikuväävli mass tiku kogumassist?
\textit{NB:} Mõõtejoonlaua kasutamine on keelatud!

\hint

\solu
Koostame kahest tikust kangkaalu. Asetame ühe tiku teise peale risti. Alumist tik-

ku hoiame nii, et terav serv oleks ülespoole, mille peale toetub ülemine tikk (vt.

joonist). Määrame ülemise tiku tasakaalupunkti ning märgime selle ära.

$l_1$ $l_2$

Lahenduses eeldame, et tikuväävli raskuskese asub tikuväävli keskpunktis ning see ühtib tiku puidust osa otspunktiga. Mõõdame tiku osade pikkused kummalgi pool tasakaalupunkti. Seal pool, kus asub tikuväävel, mõõdame kuni tikuväävli keskpunktini. Tähistame selle osa pikkuse $l_1$ ning teise osa pikkuse $l_2$. Mõõtmiseks märgime paberile tiku pikkuse ning tasakaalupunkti ning vaatame mitu tiku laiust kumbki osa pikk on. Asetame ühe tiku märgitud punkti ning tõstame teise tiku tihedalt tema kõrvale jne.

Tikuväävlile mõjub raskusjõud, mille rakenduspunkt asub toetuspunktist kaugusel $l_1$ ning tiku puidust osale raskusjõud, mille rakenduspunkt asub toetuspunktist kaugusel lx (tiku puidust osa masskese). Tähistades tikuväävli mass m ning tiku puidust osa mass M , saame kirjutada kangi tasakaalu seose kujul

mgl1 = M glx, kus lx = $l_1$ + $l_2$ $-$ $l_1$ = $l_2$ $-$ $l_1$ 2 2

Et leida, mitu protsenti moodustab tikuväävli mass tiku kogumassist, jagame tikuväävli massi m tiku kogumassiga M + m:

M + m = M + M lx = M ($l_1$ + lx) $l_1$ $l_1$

seega,

p= m = ml1

m + M M ($l_1$ + lx)

Asendades masside suhte m/M jõuõlgade suhtega lx/$l_1$, saame

p = lx $l_1$ + lx

Asendades viimasest avaldisest lx, saame

p = $l_2$ $-$ $l_1$ $l_2$ + $l_1$

Tiku osade pikkused suhtuvad keskmise tiku korral $l_1$ = 8 tiku laiust ning $l_2$ = 12 tiku laiust. Tulemus sõltuvalt tikust on 15 \% $-$ 25 \%.

\textit{Žürii hindamisskeem on kogumikus lk 85; siia ei ole seda ümber kirjutatud.}
\probend
